\chapter{Imperial China and the Tributary World Order}

\section{The dynastic cycle}

Traditional Chinese historians told history as a sequence of dynasties, each founded by a
virtuous ruler who received the Mandate of Heaven and each ending when a corrupt ruler lost
it. Modern historians treat this ``dynastic cycle'' as a \emph{narrative model}, not a law,
but it was how the actors themselves understood politics.

\begin{center}
\begin{tikzpicture}[st/.style={draw, thick, rounded corners, align=center, font=\small, text width=3.1cm, fill=yellow!12}]
  \node[st] (a) at (90:3) {new dynasty\\receives the Mandate};
  \node[st] (b) at (18:3.4) {strong rule:\\land reform, public works};
  \node[st] (c) at (-54:3.2) {decline: corruption,\\taxes, neglect};
  \node[st] (d) at (-126:3.2) {disasters, famine,\\rebellion, invasion};
  \node[st] (e) at (162:3.4) {Mandate lost;\\period of disorder};
  \draw[->, very thick] (a) to[bend left=15] (b);
  \draw[->, very thick] (b) to[bend left=15] (c);
  \draw[->, very thick] (c) to[bend left=15] (d);
  \draw[->, very thick] (d) to[bend left=15] (e);
  \draw[->, very thick] (e) to[bend left=15] (a);
  \node[font=\small\itshape] at (0,0) {\zht{天命}};
\end{tikzpicture}
\end{center}

\begin{important}[Why the cycle mattered]
\begin{itemize}
  \item It made rebellion legitimate \emph{if it succeeded}: success proved Heaven's
        approval (``the winner is king, the loser is bandit'').
  \item Natural disasters (floods, earthquakes, comets) were read as signs of Heaven's
        displeasure, so they had political meaning.
  \item Even conquest dynasties (Mongol Yuan, Manchu Qing) claimed the Mandate.
\end{itemize}
\end{important}

\section{Dynasties overview}

Know the order, rough dates and one or two defining features of each major dynasty.
(Mnemonic, to the tune of ``Frère Jacques'': ``Shang, Zhou, Qin, Han; Shang, Zhou, Qin,
Han; Sui, Tang, Song; Sui, Tang, Song; Yuan, Ming, Qing, Republic\ldots'')

\begin{center}
\begin{tabular}{llp{9cm}}
\toprule
Dynasty & Dates & Key features \\
\midrule
Shang \zht{商} & c.\ 1600--1046 BCE & bronze, oracle bones, ancestor cult \\
Zhou \zht{周} & 1046--256 BCE & Mandate of Heaven; ``feudal'' order collapses into Spring and Autumn and Warring States; Confucius, Laozi \\
Qin \zht{秦} & 221--206 BCE & first unification; First Emperor; Legalism; standardization; Great Wall linked \\
Han \zht{漢} & 206 BCE--220 CE & Confucian state; Silk Road; paper; ``Han'' becomes the name of the Chinese people \\
Six Dynasties & 220--589 & division; nomadic states in the north; Buddhism spreads \\
Sui \zht{隋} & 581--618 & reunification; Grand Canal; first exams \\
Tang \zht{唐} & 618--907 & cosmopolitan empire; capital Chang'an; golden age of poetry; model for Korea and Japan \\
Song \zht{宋} & 960--1279 & commercial revolution; printing; gunpowder; Neo-Confucianism; exam-based officialdom \\
Yuan \zht{元} & 1271--1368 & Mongol rule (Khubilai Khan); Marco Polo; attempted invasions of Japan \\
Ming \zht{明} & 1368--1644 & Han Chinese restoration; Forbidden City; Zheng He's voyages; Great Wall rebuilt \\
Qing \zht{清} & 1644--1912 & Manchu rule; largest territorial extent; Western contact and decline \\
\bottomrule
\end{tabular}
\end{center}

\subsection{The Qin unification}

The state of Qin in the far west adopted Legalist reforms (Shang Yang, 4th c.\ BCE):
peasants organized into mutual-responsibility groups, ranks awarded for military merit,
harsh punishments. In 221 BCE King Zheng conquered the other six states and took a new
title, \emph{Huangdi} (\zht{皇帝}, ``august emperor''), proclaiming himself the First
Emperor (Qin Shi Huang). The title lasted until 1912. The Terracotta Army, found in 1974,
guards his tomb near Xi'an.

\subsection{The Tang as the East Asian model}

The Tang (618--907) is the period when the shared East Asian toolkit was most actively
exported:

\begin{itemize}
  \item a codified legal system (the Tang Code, 624/653), copied in Japan's
        \emph{ritsuryō} codes and in Goryeo law;
  \item the grid-plan capital Chang'an, copied at Nara (\ja{平城京}) and Kyoto
        (\ja{平安京}) in Japan and Gyeongju in Silla;
  \item Buddhist institutions, poetry, music, dress and court ranks;
  \item foreign students: Japanese embassies (\ja{遣唐使}, kentōshi) and Silla students
        studying in Chang'an, some passing the Tang exams.
\end{itemize}

\subsection{The Song commercial revolution}

The Song is sometimes called ``early modern'': cities of over a million, paper money,
printing of books on a large scale, improved rice strains from Champa, the compass, and an
examination elite that displaced the old aristocracy. Northern territory was lost to
Jurchen invaders in 1127, and the dynasty moved south to Hangzhou.

\section{Conquest dynasties and the steppe}

The fundamental geopolitical problem of imperial China was the frontier with the steppe
nomads. Policy alternated between war, walls, marriage alliances and payment.

\begin{center}
\begin{tabular}{lll}
\toprule
Steppe/forest people & Dynasty founded in China & Period \\
\midrule
Xianbei & Northern Wei & 386--534 \\
Khitan & Liao & 916--1125 \\
Jurchen & Jin & 1115--1234 \\
Mongols & Yuan & 1271--1368 \\
Manchus (Jurchen descendants) & Qing & 1644--1912 \\
\bottomrule
\end{tabular}
\end{center}

\begin{remark}
The ``New Qing History'' school argues that the Qing were not simply sinicized; they ruled
as Confucian emperors to the Chinese, as Khans to Mongols and as patrons of Tibetan
Buddhism to Tibetans. This matters for current debates over Tibet and Xinjiang.
\end{remark}

\section{The tributary system}

\begin{definition}[Tributary system]
The \emph{tributary system} is the Western scholarly name (from John K.\ Fairbank, 1940s)
for the way imperial China organized relations with foreign states. Foreign rulers sent
embassies bearing ``tribute'' (local products) and performed the kowtow (\zht{叩頭}, three
kneelings and nine head-knockings) to the emperor. In return the emperor granted
\emph{investiture} (recognition as king), a calendar, seals, and lavish gifts, and allowed
trade.
\end{definition}

\begin{center}
\begin{tikzpicture}
  \draw[thick, fill=orange!35] (0,0) circle (1.3);
  \draw[thick] (0,0) circle (3.2);
  \draw[thick] (0,0) circle (4.9);
  \node[align=center, font=\small] at (0,0) {\textbf{Emperor}\\\zht{天子}\\``Son of\\Heaven''};
  \node[font=\small, align=center] at (0,2.05) {Sinic zone:\\Korea, Vietnam, Ryukyu,\\(Japan at times)};
  \node[font=\small, align=center] at (0,-2.15) {regular tribute,\\shared script and classics};
  \node[font=\small, align=center] at (0,4.0) {Inner Asian zone:\\Mongols, Tibetans};
  \node[font=\small, align=center] at (0,-4.0) {Outer zone:\\SE Asia, Europeans};
  \node[font=\scriptsize, align=left, anchor=west] at (5.2,1) {civilization (\zht{華})};
  \node[font=\scriptsize, align=left, anchor=west] at (5.2,-1) {``barbarians'' (\zht{夷})};
  \draw[->, thick] (5.2,0.75) -- (5.2,-0.75);
\end{tikzpicture}
\end{center}

\subsection{How it worked in practice}

\begin{itemize}
  \item \textbf{Korea} (Joseon) was the model tributary: embassies to Beijing several times a
        year, kings invested by the Ming/Qing emperor. Joseon elites called it
        \emph{sadae} (\ko{사대}, \zht{事大}, ``serving the great''). After the Ming fell,
        many Joseon literati privately regarded Korea as the true heir of Ming civilization
        (``Little China'', \zht{小中華}) and the Manchus as barbarians.
  \item \textbf{Vietnam} paid tribute to China while calling its own ruler ``emperor''
        at home and receiving tribute from smaller neighbours: a mini tributary system.
  \item \textbf{Ryukyu} (Okinawa) paid tribute to China \emph{and}, after 1609, was
        controlled by the Japanese domain of Satsuma, a ``dual subordination'' hidden from
        Chinese envoys.
  \item \textbf{Japan} mostly stayed outside. The Ashikaga shogun Yoshimitsu accepted
        investiture as ``King of Japan'' (1402) to gain trading rights, which later
        Japanese historians saw as humiliating. Tokugawa Japan traded with China through
        Nagasaki without formal tribute, and built its own Japan-centred order.
\end{itemize}

\begin{important}[Debates]
\begin{itemize}
  \item Was it really about trade? Tribute missions were often profitable for the
        foreign side, since imperial gifts exceeded the tribute.
  \item Was it a ``system'' at all? Critics argue Fairbank projected a neat model on
        varied, pragmatic relationships, and that Qing practice differed on each frontier.
  \item David Kang argues that the hierarchical order produced relatively long periods of
        peace among the Sinic states (c.\ 1368--1841), in contrast to the European
        balance-of-power system.
\end{itemize}
\end{important}

\subsection{Zheng He}

Between 1405 and 1433 the Ming eunuch admiral Zheng He led seven enormous fleets to
Southeast Asia, India, Arabia and East Africa, collecting tribute and exotic gifts (a
giraffe was presented to the emperor as a mythical \emph{qilin}). The voyages were stopped
for fiscal and political reasons. They are often contrasted with Portuguese voyages later
in the century, and are now used in PRC diplomacy as a symbol of ``peaceful rise''.

\section{Comparative governance}

\begin{center}
\begin{tabular}{lp{3.6cm}p{3.6cm}p{3.6cm}}
\toprule
 & China (Ming--Qing) & Korea (Joseon) & Japan (Tokugawa) \\
\midrule
Ruler & emperor, absolute in theory & king, constrained by yangban officials & emperor symbolic; shogun rules \\
Elite & exam-selected scholar-officials & hereditary yangban who also sit exams & hereditary samurai \\
Ideology & Neo-Confucian orthodoxy & strict Neo-Confucian orthodoxy & Neo-Confucianism plus Shinto, Buddhism \\
Structure & centralized bureaucracy & centralized bureaucracy & ``centralized feudalism'' (shogun + domains) \\
\bottomrule
\end{tabular}
\end{center}

\section*{Summary}
\begin{itemize}
  \item The dynastic cycle and the Mandate of Heaven framed Chinese political legitimacy.
  \item Know the dynasty sequence; Qin unified, Han Confucianized, Tang was the regional
        model, Song commercialized, Yuan and Qing were conquest dynasties.
  \item The steppe frontier was China's central security problem.
  \item The tributary system was a hierarchical, China-centred diplomatic order: Korea was
        its model member, Japan largely stayed outside.
  \item Scholars debate whether ``tributary system'' describes reality or a later model.
\end{itemize}
